\section{Validity and limits}
\label{sec:limits}

\paragraph{Fluent runs.}
\begin{itemize}
\item \textbf{Convergence}: $(\max C_D - \min C_D)/\overline{C_D}$ of \zone{wall\_nose} over the last 500
iterations is \val{conv-lo}--\val{conv-hi}\,\% in all eight runs, below the \SI{0.1}{\percent} of REPORT-4.
\item \textbf{Wall resolution}: largest $y^+$ \val{yplus-nose} on the nose and \val{yplus-coupler} on the
coupler, so the boundary layer is resolved down to the wall as SST requires.
\item \textbf{Mesh}: the REPORT-4 recipe, 174\,380 (A) and 174\,155 (B) cells, maximum skewness 0.90 and
0.94. No mesh study was made for these two noses; the recipe was checked on study~A only.
\item \textbf{Tip sawtooth} in $C_p$, $q_w$ and $h$ between 3 and \SI{15}{mm}, at the edge of the tip
refinement. Smoothed in the delivered tables; a \SI{40}{mm} refinement sphere would remove it. Integrated
forces and the stagnation values are not affected.
\item \textbf{Smooth wall}. The flight model uses a \SI{20}{\micro m} roughness for the friction.
\end{itemize}

\paragraph{From three runs to the whole flight.}
\begin{itemize}
\item \textbf{Zero incidence only} (2D axisymmetric). The normal force comes from slender-body theory,
which is a first estimate for a nose alone; a 3D run at incidence would be needed to refine it.
\item \textbf{Transonic rise} between $M = 0.95$ and 1.1 is not resolved. It affects the drag for about
\SI{1}{s} and none of the load cases.
\item \textbf{Hot day} is not run in Fluent: the nominal $C_p(x, M)$ is applied with the hot-day $p_\infty$,
$q$ and $M$. Above $M = \val{pt-M147-M}$ $C_p$ is held, which puts the hot-day tip pressure
\SI{2}{\percent} low (Section~\ref{sec:pressure}).
\item \textbf{Inside pressure} of the sealed nose is bounded, not computed: between the pad temperature
and the mean inner-face temperature. Leaks or vents would bring it towards $p_\infty$ and reduce
the outward loads.
\end{itemize}

\paragraph{Heating.}
\begin{itemize}
\item Fluent gives the heating at burnout only, steady. Along the flight the calibrated correlations are
used; they match Fluent at burnout within a few per cent but are not checked at other instants.
\item Wall models: one-dimensional conduction through the shell (no conduction along it, adiabatic
inside, no radiation), which is conservative for the shell peaks; and a single temperature for the tip
insert, which is its mean temperature: the front of the insert runs hotter than the value given.
\item Material data are assumptions of report~1 and must be confirmed by ST.
\end{itemize}

\paragraph{Flight model.} The trajectory is the one of report~2 for the 720\,mm noses, with the drag of
the rest of the vehicle scaled so that the 650\,mm baseline reaches \SI{9}{km}. A change of the
vehicle mass or motor changes $q$, and the pressures and forces scale with it.
